%% hopfion-framework/strong/colour.tex
%% Sources: main_paper5.tex
%% Colour from CP^2 Hopf fibration, Z_3 centre, N_c=3, confinement
%% Auto-generated from project files. Edit source papers to update.
%%
%% Status tags: \status{published}, \status{open}, \status{conjectured}
%% \source{PaperN, label} — where the proof lives
%% \residual{X\%} or \residual{exact} — numerical accuracy


%════════════════════════════════════════════════════════════════════
% Colour Structure — Paper V
%════════════════════════════════════════════════════════════════════

\begin{remark}[Parallel with the electroweak sector]
\label{P5:rem:ew_colour_parallel}
For the electroweak sector, $U(1)_Y\cap 2I = Z_2$ (the centre of the
binary icosahedral group) was the key intersection that pinned the
hypercharge coupling weight $w_Y=1$ at $C_5$ (Section~5.1.2 of~\cite{Paper1}).
For the colour sector, the analogous intersection is
$\mathrm{SU}(3)_C\cap 2I = Z_3$ (the $\mathbb{Z}_3$ centre of
$\mathrm{SU}(3)$), which pins the colour coupling weight at the
same condensate generator.
\end{remark}
\status{published}
\source{P5:rem:ew_colour_parallel}

\begin{proposition}[Colour level shift]
\label{P5:prop:colour_shift}
Every quark receives a universal half-integer shift
\begin{equation}
  \delta n^{(C)} = \frac{1}{2}
  \label{P16:eq:dncol_exact}
\end{equation}
relative to the lepton at the same Hopf tower level.
\end{proposition}
\status{published}
\source{P5:prop:colour_shift}

\begin{remark}[Comparison with the old formula]
\label{P5:rem:old_formula_comparison}
The previous estimate $\delta n^{(C)} = \tfrac{2}{3}\cdot\tfrac{\deff-2}{2}$
used $\deff = 3.5$ and coincidentally gave $1/2$, but with
$\deff^{\mathrm{exact}}=43/14$ gives the wrong result $5/14$.
Proposition~\ref{P5:prop:colour_shift} gives the exact derivation that
does not pass through $\deff$.
\end{remark}
\status{published}
\source{P5:rem:old_formula_comparison}

\begin{equation}
  \LQCD \approx \LUV\,\exp\!\left(-\frac{\vp^{43/7}}{b_0 C_F \sin(\pi/5)}\right),
  \label{P16:eq:LQCD_holonomy}
\end{equation}

%════════════════════════════════════════════════════════════════════
% Colour Exclusion Theorem — Paper XVI
%════════════════════════════════════════════════════════════════════

\begin{theorem}[Colour exclusion in the $\QH=2$ sector]
\label{P16:thm:colour_exclusion}
The colour $\mathbb{Z}_3$ automorphism $\sigma$ does not extend to any
automorphism of $\twoI$. Consequently, the colour charge is not a
well-defined quantum number for any field configuration in the
$\QH=2$ sector of the density-feedback Hopfion, and every $\QH=2$
configuration is a colour singlet.
\end{theorem}
\status{published}
\source{P16:thm:colour_exclusion}

\begin{remark}[Dynkin-diagram formulation]
\label{P16:rem:dynkin_colour}
The colour $\mathbb{Z}_3$ of $\twoT$ is the order-$3$ graph automorphism
of $\widehat{E}_6$ (the three equal arms rotate into each other).
The automorphism group of $\widehat{E}_8$ (McKay graph of $\twoI$) is
trivial: the $\widehat{E}_8$ diagram is a linear chain with a single
branch, and no non-identity node permutation preserves all adjacencies.
Therefore $\mathrm{Aut}(\widehat{E}_8)=\{1\}$, confirming by a different
route that the colour $\mathbb{Z}_3\subset\mathrm{Aut}(\widehat{E}_6)$
has no counterpart at the $\widehat{E}_8$ level.
\end{remark}
\status{published}
\source{P16:rem:dynkin_colour}

\begin{remark}[Colour emergence in the $\QH=2\to\QH=3$ transition]
\label{P16:rem:symmetry_breaking_colour}
Theorem~\ref{P16:thm:colour_exclusion} supplies the group-theoretic statement
underlying the physical picture: colour charge \emph{emerges} in the
transition $\QH=2\to\QH=3$, i.e.\ in the symmetry breaking
$\twoI\to\twoT$. Before the breaking, $\sigma$ is not a symmetry of
anything in the $\QH=2$ sector; after it, $\sigma$ is the defining
automorphism of the $\twoT$ colour structure.
\end{remark}
\status{published}
\source{P16:rem:symmetry_breaking_colour}

\begin{remark}[The $\QH=1$ sector and the complete ADE hierarchy]
\label{P16:rem:qh1_preview}
Theorem~\ref{P16:thm:colour_exclusion} concerns only the $\QH=2$ sector.
The $\QH=1$ neutrino sector (Paper XVII) has McKay group $\mathbb{Z}_2$;
colour exclusion holds there by coprimality $\gcd(2,3)=1$.
The three non-vacuum sectors complete the ADE top tier:
\[
  \underbrace{\mathbb{Z}_2\leftrightarrow\widehat{A}_1}_{\QH=1,\;\text{neutrino}}
  \qquad
  \underbrace{\twoI\leftrightarrow\widehat{E}_8}_{\QH=2,\;\text{charged lepton}}
  \qquad
  \underbrace{\twoT\leftrightarrow\widehat{E}_6}_{\QH=3,\;\text{baryon}}.
\]
\end{remark}
\status{published}
\source{P16:rem:qh1_preview}

\begin{proposition}[Topological holonomy]
\label{P5:prop:holonomy}
The correct holonomy for QCD confinement is the topological quantity
computed by Method~(iii).
For a single colour flux tube, the Dirac quantisation condition
requires $\Phi_{\mathrm{topo}} = \pi$ (one unit of colour magnetic
flux), and Method~(iii) gives $\Phi_{\mathrm{topo}}/\pi\approx1.017$,
consistent with this quantisation to $1.7\%$.
The Wilson loop for the colour flux tube is
\begin{equation}
  \langle W\rangle
  = e^{-i\Phi_{\mathrm{topo}}}
  \approx e^{-i\pi}
  = -1,
  \label{P15:eq:Wilson_loop}
\end{equation}
which is the correct $\mathbb{Z}_2$ centre value for quark confinement
(the centre-symmetry criterion for the confined phase).
\end{proposition}
\status{published}
\source{P5:prop:holonomy}

\begin{remark}[Relation to $K_{\mathrm{th}}$]
\label{P5:rem:kth_relation}
The threshold factor $K_{\mathrm{th}}=\sin(\pi/5)$ (Theorem~\ref{P5:thm:Kth})
is the WZW vacuum amplitude and modifies the \emph{coupling}, not the
holonomy.  The holonomy is topologically quantised at $\pi$; the
coupling runs with scale and receives the $K_{\mathrm{th}}$ correction
at the phase transition.  These are orthogonal corrections.
\end{remark}
\status{published}
\source{P5:rem:kth_relation}


%════════════════════════════════════════════════════════════════════
% WZW Structure of the Q_H=3 Sector — Paper XV
%════════════════════════════════════════════════════════════════════

\begin{remark}[Clarification on the McKay chain]
\label{P15:rem:mckay_chain_clarification}
The $\widehat{E}_6$ diagram (7~nodes) and the affine
$\widehat{A}_2=\widehat{\mathfrak{sl}(3)}$ diagram (3~nodes,
underlying $\SU(3)_1$) are \emph{different} diagrams.
The WZW model directly associated to $\twoT$ via the McKay
correspondence is $(E_6)_1$, not $\SU(3)_1$.
The connection to $\SU(3)_1$ is established through a
conformal embedding, which is
the $E_6\supset\SU(3)^3$ branching at level~1.
The same structural convention appears in Papers~I--XIV,
where the McKay graph of $\twoI$ is $\widehat{E}_8$ (9~nodes,
giving $(E_8)_1$ with $c=8$), while $\SU(2)_3$ ($c=9/5$)
is identified via the Pentagon number $k+2=5$ of Paper~V
rather than through a conformal embedding of $(E_8)_1$.
The $\QH=3$ sector improves on this: the bridge to $\SU(3)_1$
is an explicit conformal embedding with exact central-charge match.
\end{remark}
\status{published}
\source{P15:rem:mckay_chain_clarification}

\begin{definition}
\label{P15:def:closes_tmatrix}
An integer $Q\geq 1$ \emph{closes} the T-matrix of a WZW model
if $Q\,T_j$ is an odd multiple of $\tfrac{1}{4}$ for every primary~$j$
(equivalently, $\cos(2\pi Q T_j)=0$ for all~$j$).
\end{definition}
\status{published}
\source{P15:def:closes_tmatrix}

\begin{theorem}[T-matrix closure at $(E_6)_1$ level]
\label{P15:thm:QgroupE6}
The minimum closing integer for $(E_6)_1$ (i.e.\ the minimum $Q\geq1$
such that $Q\,T_j$ is an odd multiple of $\tfrac{1}{4}$ for every primary $j$)
is $Q_{\mathrm{group}}^{(3)}=3$.
\end{theorem}
\status{published}
\source{P15:thm:QgroupE6}
\residual{exact}

\begin{theorem}[T-matrix closure at $\SU(3)_1$ level]
\label{P15:thm:Qgroup3}
The minimum closing integer for $\SU(3)_1$ is also
$Q_{\mathrm{group}}^{(3)}=3$, independently of the $(E_6)_1$ result.
\end{theorem}
\status{published}
\source{P15:thm:Qgroup3}
\residual{exact}

\begin{remark}[Robustness of $Q_{\mathrm{group}}^{(3)}=3$]
\label{P15:rem:qgroup_robustness}
The closing integer $Q=3$ is obtained independently at both levels
of the WZW chain.
It equals $|2T|/8=3$ (where~8 is the order of the
octahedral faces).
The agreement between the $(E_6)_1$ and $\SU(3)_1$ results
confirms the internal consistency of the WZW chain.
\end{remark}
\status{published}
\source{P15:rem:qgroup_robustness}

\begin{theorem}[Both quark charges from the WZW chain]
\label{P15:thm:charge}
Both fractional quark charges emerge at different levels of the
WZW chain, with no free parameters:
\begin{enumerate}[noitemsep,label=(\roman*)]
\item $(E_6)_1$ level: $h_{\mathbf{27}}=\tfrac{2}{3}$
  gives the up-quark charge $|Q_u|=\tfrac{2}{3}$.
\item $\SU(3)_1$ level ($\SU(3)_L$ electroweak factor):
  \begin{equation}
    h_{(1,0)} = \frac{p^2+pq+q^2+3p+3q}{12}\bigg|_{p=1,q=0}
              = \frac{4}{12} = \frac{1}{3}
    \label{P15:eq:hdown}
  \end{equation}
  gives the down-quark charge $|Q_d|=\tfrac{1}{3}$.
\item These are consistent via branching:
  $h_{\mathbf{3}}+h_{\bar{\mathbf{3}}}
  =\tfrac{1}{3}+\tfrac{1}{3}=\tfrac{2}{3}=h_{\mathbf{27}}$.
\end{enumerate}
\end{theorem}
\status{published}
\source{P15:thm:charge}
\residual{exact}

\begin{remark}[Both charges, no free parameters]
\label{P15:rem:both_charges_no_params}
Neither $|Q_u|=2/3$ nor $|Q_d|=1/3$ is an input; both are outputs
of the McKay correspondence and conformal embedding.
The ratio $|Q_u|/|Q_d|=2$ follows from the $E_6\supset\SU(3)^3$
branching.
\end{remark}
\status{published}
\source{P15:rem:both_charges_no_params}

\begin{theorem}[Electric charge from the trinification generator]
\label{P15:thm:charge_generator}
Let $Q = T_{3L}+T_{8L}/\!\sqrt{3}$ be the unique electric-charge
generator of $\SU(3)_L\subset(E_6)_1$ (traceless over $\mathbf{3}_L$
and compatible with the $\mathbb{Z}_3$ centre; no free parameters).
On the three weight states of the fundamental $\mathbf{3}_L$,
\begin{equation}
  Q \in \Bigl\{\tfrac{2}{3},\,-\tfrac{1}{3},\,-\tfrac{1}{3}\Bigr\},
  \label{P15:eq:trinification_charges}
\end{equation}
the standard $(u,d,D)$ charge assignment. The coefficient vector of
$Q$ satisfies $\mu = 2\,\omega_1$ where $\omega_1$ is the highest
weight of $\mathbf{3}_L$.
\end{theorem}
\status{published}
\source{P15:thm:charge_generator}
\residual{exact}

\begin{corollary}[Quark charges equal conformal weights, derived]
\label{P15:cor:charge_eq_weight}
Both quark charges and both conformal weights are computed from the
single number $(\omega_1,\omega_1)_{A_2}=\tfrac23$:
\begin{equation}
  \boxed{|Q_d| = h_{(1,0)} = \tfrac13,
  \qquad
  |Q_u| = Q(\omega_1) = 2h_{(1,0)} = h_{\mathbf{27}} = \tfrac23.}
  \label{P15:eq:charge_eq_weight}
\end{equation}
The equality of conformal weight and electric charge is a consequence
of the trinification embedding with $A_2$ representation theory.
\end{corollary}
\status{published}
\source{P15:cor:charge_eq_weight}
\residual{exact}

\begin{remark}[Why $N=3$ is essential]
\label{P15:rem:N3_unique}
For $\SU(N)_1$, the conformal weight of the fundamental equals the
minimal $\mathbb{Z}_N$-centre charge if and only if
\begin{equation}
  \frac{N-1}{2N}=\frac{1}{N}
  \quad\Longleftrightarrow\quad
  N=3.
  \label{P15:eq:N3unique}
\end{equation}
The identification of Corollary~\ref{P15:cor:charge_eq_weight} is therefore
special to $N=3$, the value independently forced by the McKay
correspondence $\twoT\leftrightarrow\widehat{E}_6$.
\end{remark}
\status{published}
\source{P15:rem:N3_unique}

\begin{remark}[Leading-order mass degeneracy]
\label{P15:rem:leading_order_mass_degeneracy}
In $\SU(3)_1$: $T_{(1,0)}=T_{(0,1)}=\tfrac{1}{4}$ (identical
T-matrix phases); by the same logic as Papers~III--IV~\cite{Paper3,Paper4},
this predicts $m_u=m_d$ at leading order.
In $(E_6)_1$: $T_{\mathbf{27}}=T_{\overline{\mathbf{27}}}=\tfrac{5}{12}$
(same degeneracy).
The prediction $m_u\approx m_d$ is consistent with the measured
ratio $m_u/m_d\approx0.46$, the smallest quark mass ratio,
generated by subleading isospin-breaking effects.
\end{remark}
\status{published}
\source{P15:rem:leading_order_mass_degeneracy}

%════════════════════════════════════════════════════════════════════
% 2T Group Structure: Subgroup, 24-Cell, Branching — Paper XV
%════════════════════════════════════════════════════════════════════

\begin{theorem}[$\twoT\subset\twoI$]
\label{P15:thm:subgroup}
The binary tetrahedral group $\twoT$ (order~24) is a subgroup of
the binary icosahedral group $\twoI$ (order~120).
The index $[\twoI:\twoT]=5=k+2$ at $k=3$.
\end{theorem}
\status{published}
\source{P15:thm:subgroup}

\begin{remark}[Index and structure]
\label{P15:rem:index_and_structure}
The index $[\twoI:\twoT]=120/24=5$ equals the pentagonal
symmetry number $k+2\big|_{k=3}=5$ of the WZW level.
The five cosets of $\twoT$ in $\twoI$
correspond to the five lepton primaries after including the
vacuum $j=0$, and the restriction $\twoI\to\twoT$ is the algebraic
content of passing from the lepton sector ($\QH=2$) to the
quark sector ($\QH=3$).
\end{remark}
\status{published}
\source{P15:rem:index_and_structure}

\begin{proposition}[24-cell structure of $\twoT$]
\label{P15:prop:24cell}
The 24 elements of $\twoT\subset\SU(2)$, as unit quaternions,
are the vertices of the 24-cell in $\mathbb{R}^4$.
They decompose as $Q_8 = \{\pm1,\pm\mathbf{i},\pm\mathbf{j},\pm\mathbf{k}\}$
(8 elements) and the 16 Hurwitz units
$\tfrac{1}{2}(\pm1\pm\mathbf{i}\pm\mathbf{j}\pm\mathbf{k})$.
The 24-cell has edge length~$1$ and circumradius~$1$;
its edge-to-circumradius ratio is exactly~$1$
(compared to $1/\vp$ for the 600-cell of $\twoI$).
\end{proposition}
\status{published}
\source{P15:prop:24cell}

\begin{remark}[The $F_4$ root system and Coxeter number for quarks]
\label{P15:rem:F4}
The McKay correspondence gives $\twoT\leftrightarrow F_4$:
\begin{equation}
  |\twoT|=24 = \tfrac12\cdot|R(F_4)|=\tfrac12\cdot48.
  \label{P15:eq:2T_F4}
\end{equation}
The dual Coxeter number $h^\vee(F_4)=9=N_3^2$.
The ratio $h(F_4)/h(E_8)=12/30=2/5=(k+2)^{-1}$ at $k=3$, consistent
with the index $[\twoI:\twoT]=5=k+2$.
\end{remark}
\status{published}
\source{P15:rem:F4}

\begin{remark}[Integer normalisation from the 24-cell]
\label{P15:rem:integer_norm_24cell}
The 600-cell's edge-to-circumradius ratio $1/\vp$ produces the
golden-ratio arithmetic of the $\QH=2$ sector: $\vp^9\sqrt{\Ja J_4}=10$.
The 24-cell's ratio of exactly~$1$ (no $\vp$) corresponds to the
$\QH=3$ normalisation identity $\vp^9\sqrt{\Ja^{(3)}J_4^{(3)}}=3$
having an \emph{integer} right-hand side.
The $\vp$ factors in the $\QH=3$ identity come entirely from the
condensate (the $\vp^9$ prefactor), not from the quark-sector group.
This reflects the 24-cell's rational arithmetic: there is no
golden ratio in the geometry of $\twoT$.
Since the colour gauge structure $\SU(3)_1$ arises from the conformal
embedding $(E_6)_1\supset\SU(3)_1^{\times3}$,
and $(E_6)_1$ is itself the McKay dual of the $\vp$-free group $\twoT$,
the colour algebra inherits no $\vp$-dependence at the group-theoretic
level; this is consistent with $\vp$ being wholly absent from QCD's
own gauge structure (the Gell-Mann algebra, the running of $\alpha_s$,
and the colour-confinement mechanism), even though $\vp$ enters the
condensate energy scale $\Efb^{(3)}=2N_3^2/\vp^{17}$ that sets the
absolute mass scale (open problem~(iv)).
\end{remark}
\status{published}
\source{P15:rem:integer_norm_24cell}

\begin{proposition}[Branching rules $\twoI\to\twoT$, proved by character theory]
\label{P15:prop:branching}
\begin{equation}
\renewcommand{\arraystretch}{1.3}
\begin{array}{rll}
  j=0\phantom{/2}& \to & \Gamma_1 \\
  j=\tfrac12 & \to & \Gamma_2 \\
  j=1\phantom{/2} & \to & \Gamma_3 \\
  j=\tfrac32 & \to & \Gamma_{2\omega}\oplus\Gamma_{2\omega^2} \\
  j=2\phantom{/2} & \to & \Gamma_\omega\oplus\Gamma_{\omega^2}\oplus\Gamma_3 \\
  j=\tfrac52 & \to & \Gamma_2\oplus\Gamma_{2\omega}\oplus\Gamma_{2\omega^2}
\end{array}
\label{P15:eq:branching_table_P15}
\end{equation}
All multiplicities are $0$ or $1$; the restriction $\twoI\to\twoT$ is
multiplicity-free for $j\le\tfrac52$.
\end{proposition}
\status{published}
\source{P15:prop:branching}

\begin{remark}[Key physical features of the branching rules]
\label{P15:rem:branching_key_features}
Three entries of the branching table have direct physical content.

\emph{(i) $j=\tfrac12\to\Gamma_2$:} the $\QH=1$ fundamental maps to
colour~1. This grounds the bosonization identification (Paper~III
Proposition~10.6) across both sectors.

\emph{(ii) $j=\tfrac32\to\Gamma_{2\omega}\oplus\Gamma_{2\omega^2}$:}
the simple current of $\mathrm{SU}(2)_3$ decomposes into
colours~2 and~3 only, not colour~1.
The three colours together appear only in the triple fusion
$j=\tfrac12\otimes j=\tfrac12\otimes j=\tfrac12$
(which contains $j=\tfrac52\to\Gamma_2\oplus\Gamma_{2\omega}\oplus\Gamma_{2\omega^2}$,
all three colours simultaneously), confirming the three-tube structure of the trefoil.

\emph{(iii) $j=0,\tfrac12,1\to\Gamma_1,\Gamma_2,\Gamma_3$:}
the three lepton-generation levels of $\mathrm{SU}(2)_3$
each map to a \emph{distinct} irrep of $\twoT$.
For quarks, the generation label $(j_g,\Gamma_{2\omega^c})$ carries
T-matrix phase $T_{j_g}^{\mathrm{SU}(2)_3}\in\{-\tfrac{3}{40},\tfrac{3}{40},\tfrac{13}{40}\}$
with $Q_3=3$, giving inter-generation T-phase ratios that are $O(1)$ corrections;
the large mass hierarchy comes from the $\vp$-tower levels (open problem~(iv)(c)).
\end{remark}
\status{published}
\source{P15:rem:branching_key_features}

\begin{proposition}[Three colours from three fermionic irreps]
\label{P15:prop:colours}
The three $2$-dimensional fermionic irreps
$\Gamma_2,\Gamma_{2\omega},\Gamma_{2\omega^2}$ of $\twoT$,
related by the $\mathbb{Z}_3$ outer automorphism of $E_6$,
are the three quark colours:
\begin{equation}
  \twoI(\text{spinor}, \dim=2)\;\big\downarrow_{\twoT}\;
  \longrightarrow\; \Gamma_{2\omega^k},
  \qquad k\in\{0,1,2\},
  \label{P15:eq:branching_P15}
\end{equation}
with index $k$ labelling the quark colour.
\end{proposition}
\status{published}
\source{P15:prop:colours}

\begin{theorem}[Charge and mass from $\SU(3)_1$]
\label{P15:thm:massdegen}
Within the $\SU(3)_1$ WZW framework for the $\QH=3$ sector:
\begin{enumerate}[noitemsep,label=(\roman*)]
\item The fractional electric charge of the quark is
  $h_{(1,0)}=\tfrac{1}{3}$ (Theorem~\ref{P15:thm:charge}).
\item The T-matrix exponents satisfy
  $T_{(1,0)}=T_{(0,1)}=\tfrac{1}{4}$: the fundamental
  $\mathbf{3}$ and anti-fundamental $\bar{\mathbf{3}}$ have
  identical T-matrix phases.
\item At leading order in the WZW framework, $m_u=m_d$.
\end{enumerate}
\end{theorem}
\status{published}
\source{P15:thm:massdegen}

\begin{remark}[Consistency with observation]
\label{P15:rem:consistency_with_observation}
The prediction $m_u\approx m_d$ at leading order is consistent
with the measured ratio $m_u/m_d\approx 0.46$ (the smallest among six quarks).
The large hierarchies $m_c/m_u\sim10^3$, $m_t/m_u\sim10^5$
involve the $\vp$-tower levels and the quark Yukawa coupling,
deferred to Paper~XVI~\cite{Paper16} (open problem~(iv)(c)).
\end{remark}
\status{published}
\source{P15:rem:consistency_with_observation}
